\documentclass[11pt,a4paper]{article}
\usepackage[T1]{fontenc}
\usepackage[utf8]{inputenc}
\usepackage{lmodern}
\usepackage{microtype}
\usepackage[margin=25mm]{geometry}
\usepackage{amsmath,amssymb,amsthm,mathtools}
\usepackage{booktabs,tabularx,array}
\usepackage{enumitem}
\usepackage{xcolor}
\usepackage{xurl}
\usepackage{hyperref}
\usepackage{fancyhdr}
\usepackage{needspace}

\definecolor{linkblue}{RGB}{25,62,100}
\hypersetup{
  colorlinks=true,
  linkcolor=linkblue,
  citecolor=linkblue,
  urlcolor=linkblue,
  pdftitle={Congruence Carriers and a Sextic Reduction for Diophantine Consistency},
  pdfauthor={Research contribution for the ProveIt repository},
  pdfsubject={Explicit Diophantine forms of syntactic ZFC consistency}
}
\setlength{\headheight}{14pt}
\pagestyle{fancy}
\fancyhf{}
\fancyhead[L]{\small Explicit Diophantine consistency equations}
\fancyhead[R]{\small\thepage}
\renewcommand{\headrulewidth}{0.3pt}
\setlength{\emergencystretch}{2em}
\setcounter{tocdepth}{2}
\setlist[itemize]{topsep=5pt,itemsep=3pt}
\setlist[enumerate]{topsep=5pt,itemsep=3pt}
\allowdisplaybreaks[1]

\newtheorem{theorem}{Theorem}[section]
\newtheorem{lemma}[theorem]{Lemma}
\newtheorem{proposition}[theorem]{Proposition}
\newtheorem{corollary}[theorem]{Corollary}
\theoremstyle{definition}
\newtheorem{definition}[theorem]{Definition}
\newtheorem{remark}[theorem]{Remark}
\newcommand{\Z}{\mathbb Z}
\newcommand{\N}{\mathbb N}
\newcommand{\Zpos}{\mathbb Z_{>0}}
\newcommand{\ZFC}{\mathrm{ZFC}}
\newcommand{\Con}{\operatorname{Con}}
\newcommand{\ZFCRef}{\operatorname{Ref}}
\newcommand{\Proof}{\operatorname{Proof}}
\newcommand{\degw}{\operatorname{deg}_{\mathrm{wit}}}
\newcommand{\snapshot}{\texttt{9cca523d9f27ad39d52993f4065cbd7e58f8bdd4}}
\newcommand{\repofile}[2]{\href{https://github.com/VladimirReshetnikov/ProveIt/blob/9cca523d9f27ad39d52993f4065cbd7e58f8bdd4/Computability/DiophantineConsistency/#1}{#2}}
\newcommand{\formfile}[2]{\repofile{con-zfc-diophantine-forms/#1}{#2}}

\title{\textbf{Congruence Carriers and a Sextic Reduction}\\[3pt]
  \textbf{for Diophantine Consistency}\\[8pt]
  \large Thirteen positive unknowns in degree 28\\
  Fourteen positive unknowns in degree 24\\
  Thirty-four unrestricted unknowns in degree 6}
\author{Research contribution prepared for Vladimir Reshetnikov\\
  and the ProveIt repository}
\date{23 October 2026}

\begin{document}
\maketitle

\begin{abstract}
We give three fixed integer polynomials whose lack of solutions, in their
specified domains, is equivalent in ordinary ZFC mathematics to the
nonexistence of a finite syntactic ZFC refutation.  The first has thirteen
strictly positive unknowns, exact degree \(28\), and a literal circuit of
\(98\) binary arithmetic operations.  Its comparison point in the inspected
repository has degree \(34\) and \(93\) operations at the same number of
unknowns.  Retaining one product as a fourteenth positive coordinate
gives exact degree \(24\) in \(101\) operations, compared with the inspected
fourteen-variable degree-\(26\), \(92\)-operation endpoint.
The third is a sextic in thirty-four unrestricted integers,
the sum of squares of nineteen cubic-or-lower residuals; the inspected
sextic has thirty-five unknowns and twenty residuals.

For the positive construction, an existing monic Pell equation carries
the fixed congruence needed to make a source gap an independent positive
coordinate.  This changes the order in which source bounds, divisibility,
and signs become available.  For the sextic, two polynomial root changes
and a row reduction modulo a retained square equation allow seven
registers to be eliminated from the source quartic while keeping every
remaining residual cubic.  We prove the new reductions and reproduce
the signed auxiliary-index and squared-modulus completeness arguments.

The archive specifies coherent fixed coefficient tuples, literal
polynomials, exact checkers, and source provenance.  The consistency
corollaries use the named inherited native compiler and refutation
calibration.  Finite verification establishes algebra and data integrity;
the exact object-language \texttt{BProv} obligation recorded by ProveIt
remains a separate formalization task.  These are measured improvements
over a pinned comparison snapshot, with no global minimality or global
novelty claim.
\end{abstract}

\clearpage
\tableofcontents
\newpage

\section{The target and the three completed endpoints}
\label{sec:target}

\subsection{A fixed syntactic consistency statement}

Fix the effective proof calculus for ZFC used by the inherited refutation
calibration.  Write
\(\Proof_{\ZFC}(p,\ulcorner\bot\urcorner)\) for its finite proof-checking
relation: \(p\) codes a syntactic derivation of contradiction from ZFC
axioms.  A proof code includes the finite formulas and instances of the
axiom schemata required to check that derivation.  Define
\begin{equation}
 \ZFCRef_{\ZFC}\;:\!\Longleftrightarrow\;
 \exists p\in\N\;\Proof_{\ZFC}(p,\ulcorner\bot\urcorner),
 \qquad
 \Con(\ZFC)\;:\!\Longleftrightarrow\;\neg\ZFCRef_{\ZFC}.
 \label{eq:syntactic-con}
\end{equation}
Thus the target is syntactic consistency.  It does not refer to a
transitive model, correctness about all arithmetical statements, or an
unbounded reflection principle.  The source calibration and its proof
convention are part of the specification, not free choices made after
examining a putative solution; see
\cite{fixed-source,native-calibration,object-zfc}.

A \emph{fixed} polynomial here has integer coefficients given by finite
recipes and by authenticated numeral files.  The only quantified
parameters are the variables in its displayed witness list.  Names for
subexpressions are arithmetic registers, not additional existential
variables.  Powers occurring in a recipe for a fixed coefficient have
fixed finite exponents.

Let \(P_{13}\) be the polynomial of
Section~\ref{sec:positive-coordinate-carrier}, bound to the entire compact
source-44 tuple, and let \(P_{14}\) be its product-register lift in
Section~\ref{sec:positive14-degree24}, with the same entire tuple.
Let \(Q_{34}\) be the polynomial of
Section~\ref{sec:sextic34}, bound to the entire source-34 tuple.  In
particular, the subscripts count unknowns, and the domains are different.

\begin{theorem}[The completed polynomials]
\label{thm:main}
The following ordinary mathematical reductions hold.
\begin{enumerate}
 \item \(P_{13}\) has thirteen strictly positive unknowns, exact total
 degree \(28\), and \(98\) live binary arithmetic gates.  Its positive
 zeros give normalized native-source witnesses; conversely a finite
 refutation, through the canonical native-source construction, gives a
 positive zero of \(P_{13}\).
 \item \(P_{14}\) has fourteen strictly positive unknowns, exact total
 degree \(24\), and \(101\) live binary arithmetic gates.  Its zero set
 is the graph of the uniquely determined positive coordinate \(E=XY\)
 over the positive zero set of \(P_{13}\).
 \item \(Q_{34}\) has thirty-four unrestricted integer unknowns and exact
 degree \(6\).  Its nineteen residuals have degree at most three.  Its
 complete integer zero set is in explicit polynomial bijection with
 the zero set of the pinned forty-one-integer source quartic.
\end{enumerate}
Consequently, using the inherited source theorems identified in
Section~\ref{sec:proof-status},
\begin{align}
 \exists\mathbf x\in\Zpos^{13}\;P_{13}(\mathbf x)=0
 &\quad\Longleftrightarrow\quad\ZFCRef_{\ZFC},
 \label{eq:main-pos-existence}\\
 \exists\mathbf t\in\Zpos^{14}\;P_{14}(\mathbf t)=0
 &\quad\Longleftrightarrow\quad\ZFCRef_{\ZFC},
 \label{eq:main14-existence}\\
 \exists\mathbf z\in\Z^{34}\;Q_{34}(\mathbf z)=0
 &\quad\Longleftrightarrow\quad\ZFCRef_{\ZFC}.
 \label{eq:main-int-existence}
\end{align}
\end{theorem}

\begin{proof}
The first reduction is proved by congruence recovery, Pell-index recovery,
and constructive completeness in
Section~\ref{sec:positive-coordinate-carrier}.  The elementary auxiliary
lemmas used there are proved in Section~\ref{sec:signed-index}.
The product-register lift is proved in
Section~\ref{sec:positive14-degree24}.
The sextic reduction follows from the two inverse coordinate maps and
the retained-row identity in Section~\ref{sec:sextic34}.  Composing with
the pinned native-source and finite-refutation equivalences gives
\eqref{eq:main-pos-existence}--\eqref{eq:main-int-existence}.
\end{proof}

\Needspace{9\baselineskip}
\begin{corollary}[Ordinary ZFC-provable consistency equivalences]
\label{cor:zfc}
With the fixed proof convention and coefficients just specified, the
ordinary proofs, composed with the identified source proofs, establish
\begin{align}
 \ZFC\ \vdash\
 \left(\forall\mathbf x\in\Zpos^{13}\;
      P_{13}(\mathbf x)\ne0\right)
 &\ \longleftrightarrow\ \Con(\ZFC),
 \label{eq:zfc-pos}\\
 \ZFC\ \vdash\
 \left(\forall\mathbf t\in\Zpos^{14}\;
      P_{14}(\mathbf t)\ne0\right)
 &\ \longleftrightarrow\ \Con(\ZFC),
 \label{eq:zfc-pos14}\\
 \ZFC\ \vdash\
 \left(\forall\mathbf z\in\Z^{34}\;
      Q_{34}(\mathbf z)\ne0\right)
 &\ \longleftrightarrow\ \Con(\ZFC).
 \label{eq:zfc-int}
\end{align}
\end{corollary}

\begin{proof}
Negate the existence equivalences of Theorem~\ref{thm:main} and use
\eqref{eq:syntactic-con}.  The coordinate manipulations, divisibility
arguments, integer recurrences, and finite syntactic constructions are
ordinary mathematics available in ZFC.  This composition uses the
unbounded arithmetic source proofs, not an inference from finite
numerical tests or merely from an external truth statement.
\end{proof}

The turnstiles in Corollary~\ref{cor:zfc} state the ordinary mathematical
provability claim.  They do not assert that this archive contains an
exported derivation in ProveIt's particular \texttt{SetTheory.Form}
calculus.  Section~\ref{sec:object-language} records the exact distinction
and the still-open formal endpoint.  Neither direction of the
equivalence assumes that ZFC is consistent.

\subsection{What is being optimized}

The comparison repository is pinned at
\begin{center}
 \snapshot.
\end{center}
The following table reports the relevant inspected endpoints.  It is a
comparison of specific artifacts, not a classification of all possible
Diophantine representations.  The positive predecessor is package~61's
principal thirteen-variable equation, not its fourteen-variable
intermediate; the integer sextic predecessor is package~62
\cite{source61,source62}.

\begin{table}[htbp]
\centering
\small
\begin{tabularx}{\textwidth}{@{}lrrrX@{}}
\toprule
Domain & Unknowns & Degree & Gates & Endpoint or role\\
\midrule
\(\Zpos\) & 13 & 34 & 93 & Package 61; positive comparison point\\
\(\Zpos\) & \textbf{13} & \textbf{28} & \textbf{98}
 & This report, \(P_{13}\)\\
\(\Zpos\) & 14 & 26 & 92 & Package 65; a distinct tradeoff\\
\(\Zpos\) & \textbf{14} & \textbf{24} & \textbf{101}
 & This report, \(P_{14}\)\\
\(\Zpos\) & 15 & 24 & 89 & Package 60; a distinct tradeoff\\
\addlinespace
\(\Z\) & 41 & 4 & --- & Package 64; quartic source\\
\(\Z\) & 35 & 6 & --- & Package 62; sextic comparison point\\
\(\Z\) & \textbf{34} & \textbf{6} & --- & This report, \(Q_{34}\)\\
\(\Z\) & 30 & 8 & --- & Package 62; a distinct tradeoff\\
\bottomrule
\end{tabularx}
\caption{Measured local improvements and nearby inspected alternatives.
A dash means that no optimized gate claim is being compared.  Source
references: \cite{source61,source62,source64,source65,source60}.}
\label{tab:comparison}
\end{table}

For thirteen positive unknowns the degree falls by \(6/34\), approximately
\(17.65\%\), and the circuit gains five gates.  For degree six over
\(\Z\), one unknown and one residual are removed from the comparison
endpoint; the variable reduction is \(1/35\), approximately \(2.86\%\).
The fourteen-coordinate lift lowers degree \(26\) to \(24\) with nine
additional gates.  Relative to the fifteen-positive degree-\(24\),
\(89\)-gate equation, it removes one unknown at a cost of twelve gates.
The cheaper positive alternatives and the thirty-integer octic remain
distinct tradeoffs.

Every binary addition, subtraction, or multiplication is charged as one
gate, including multiplication by a fixed scalar and squaring.
Previously computed values may be reused.  Fixed numeral creation is
free in this arithmetic-circuit model.  It is therefore a different
measurement from bit complexity, coefficient height, expanded term
count, or the length of a source-language expression.
For \(Q_{34}\) the report claims exact variables, degree, residual count,
and expanded term count, without a gate optimum.

The pinned collection already contains equations with fewer unknowns at
much higher degrees.  These endpoints do not improve those
fewest-variable bounds and do not minimize every measure simultaneously.
The contribution is the completed arithmetic transformations and their
measured effect on the named comparison points.

\section{Two general algebraic mechanisms}
\label{sec:mechanisms}

\subsection{A fixed congruence inside a monic root}

An equation which already requires an integral root can enforce an
additional fixed congruence without a separate quotient variable.  The
following elementary observation states the exact integrality mechanism.

\begin{lemma}[Monic fixed-congruence carrier]
\label{lem:monic-carrier}
Let \(m>0\) and \(r\) be fixed integers.  Let \(H(\mathbf x)\) be an
integer polynomial, and let
\[
 F(T;\mathbf x)=T^d+
     a_{d-1}(\mathbf x)T^{d-1}+\cdots+a_0(\mathbf x)
 \in\Z[\mathbf x,T]
\]
be monic in \(T\).  Define the integer polynomial
\[
 G(t;\mathbf x)
  =m^dF\!\left(t+\frac{H(\mathbf x)-r}{m};\mathbf x\right).
\]
At every integral zero of \(G\), one has
\[
 H(\mathbf x)\equiv r\pmod m,
 \qquad
 T=t+\frac{H(\mathbf x)-r}{m}\in\Z,
 \qquad F(T;\mathbf x)=0.
\]
Conversely, an integral zero of \(F\) satisfying that congruence gives
an integral zero of \(G\) by
\(t=T-(H(\mathbf x)-r)/m\).
For strictly positive coordinates this converse also requires \(t>0\).
\end{lemma}

\begin{proof}
Multiplication by \(m^d\) clears all denominators, so \(G\) is an integer
polynomial.  At a zero, the displayed rational value of \(T\) is a root
of a monic polynomial with integer coefficients.  If \(T=a/b\) in
lowest terms with \(b>0\), multiplying \(F(a/b;\mathbf x)=0\) by \(b^d\)
shows that \(b\mid a^d\); hence \(b=1\).  Since \(t\in\Z\),
\((H(\mathbf x)-r)/m\in\Z\), which is the required congruence.
The converse is substitution, with the indicated domain condition.
\end{proof}

The simplest instance is a square equation:
\begin{equation}
 (mt+H-r)^2=m^2W(\mathbf x).
 \label{eq:monic-square-carrier}
\end{equation}
Its quotient is a rational square root of an integer and is therefore an
integer.  This conclusion does not require coprimality of \(m\) with
\(H-r\), nor factoring \(W\).

For \(P_{13}\), the already needed strong root is written
\[
 N_f=3\beta f_0+q-\beta-1,\qquad f=N_f/(3\beta),
\]
and the strong residual is
\[
 N_f^2-(3\beta)^2\{1+\Delta(ic^2)^2\}=0.
\]
Lemma~\ref{lem:monic-carrier} yields
\[
 q=\beta(3j+1)+1,\qquad J=(q-1)/\beta=3j+1.
\]
The supplied coordinates can consequently be \(F,\alpha,C\), with
\(q=F+\alpha+C\), rather than \(F,\alpha,J\) with a gap whose sign still
needs to be proved.

This is useful because \(C>0\) immediately gives the source packing bounds
once the congruence has been recovered.  The first carrier no longer
needs to establish the transport quotient's sign before those bounds
can be used.  The canonical source has \(J\equiv1\pmod3\) and a strong
root large enough that \(f_0=f-(J-1)/3>0\), which supplies the converse
domain condition.  The general lemma alone would not supply that
positivity or the native source theorem.

\subsection{Reduce a row modulo equations which are retained}

The next mechanism concerns complete systems of equations.  Degree can
fall when a residual is replaced by an equivalent representative modulo
other residuals that remain in the system.

\begin{lemma}[Retained-row degree reduction]
\label{lem:retained-row}
Let \(R_1,\ldots,R_s,G_1,\ldots,G_t\) be integer polynomials, and choose
integer polynomials \(H_{ij}\).  Put
\[
 G'_i=G_i+\sum_{j=1}^{s}H_{ij}R_j
 \qquad(1\le i\le t).
\]
The systems
\[
 R_1=\cdots=R_s=G_1=\cdots=G_t=0
 \quad\hbox{and}\quad
 R_1=\cdots=R_s=G'_1=\cdots=G'_t=0
\]
have the same zero set over \(\Z\), indeed over every commutative ring.
If all retained \(R_j\) and replaced \(G'_i\) have degree at most \(d\),
their sum of squares is an integer polynomial of degree at most \(2d\)
whose integer zeros are exactly that common zero set.
\end{lemma}

\begin{proof}
On the locus \(R_1=\cdots=R_s=0\), each \(G'_i\) equals \(G_i\).
Equivalently, the change of rows is a triangular polynomial matrix
with diagonal entries one and inverse obtained by replacing
\(H_{ij}\) with \(-H_{ij}\).  The sum-of-squares assertion follows
because a sum of integer squares is zero exactly when all its
summands are zero.
\end{proof}

Only the last sum-of-squares assertion uses the ordered integer domain;
the row-equivalence assertion does not.  The transformed sum of squares
need not equal the original polynomial.  It has the same integer zero
set, which is the property required here.

A concrete saving is
\begin{align}
 R&=F_2-f^2, \notag\\
 G&=(fo-c)^2+g-y^2, \notag\\
 G'&=F_2o^2-2foc+c^2+g-y^2,
 &G'-G&=o^2R.
 \label{eq:retained-square-example}
\end{align}
The original \(G\) is quartic in the supplied variables; \(G'\) is cubic.
The quadratic definition \(R=0\) is therefore valuable to retain, even
while other definitions are eliminated.

\subsection{Coordinate changes and the choice of registers}

A defined register \(z=E(\mathbf x)\) can be deleted by substituting
\(E(\mathbf x)\), with inverse restoration \(z=E(\mathbf x)\).
This gives a bijection of zero sets, but the substituted degree may be
larger.  A polynomial shift of one root,
\[
 t_{\mathrm{old}}=t+B(\mathbf x),
\]
is also reversible over \(\Z\) when \(B\) does not involve \(t\).
Its purpose here is to cancel a known leading part of a norm before
deleting a register.

For the integer sextic, the two useful identities are
\begin{align*}
 (t+B)^2-(k+B)B-1
   &=t^2+(2t-k)B-1,\\
 (D+ac)^2-(a^2+4a+3)c^2-1
   &=D^2+2acD-(4a+3)c^2-1.
\end{align*}
The second permits deletion of the discriminant register after changing
the main root.  Retaining \(F_2\), and using
\eqref{eq:retained-square-example}, then avoids a quartic auxiliary
row.  Section~\ref{sec:sextic34} supplies every restored coordinate and
every retained equation.

The converse degree tradeoff is also useful.  In \(P_{13}\), the
abbreviation \(E=XY\) occurs inside expressions which are subsequently
squared.  Supplying \(E\) as one more positive coordinate and retaining
the equation \(E-XY=0\) lowers those degrees.  Each positive old zero
has exactly one new coordinate \(E=XY>0\); conversely the retained
equation restores the old expression.  This graph construction produces
\(P_{14}\) without an additional semantic assumption, as proved in
Section~\ref{sec:positive14-degree24}.

Monicity, polynomial coordinate changes, and elementary row operations
are established algebraic tools.  The specific contribution here is
their placement in these arithmetic systems, including the domain and
completeness proofs which make the resulting variable and degree
counts valid.  The report does not claim that these general algebraic
facts were previously unknown.

\section{Fixed coefficients and the native source interfaces}
\label{sec:coefficients}

\subsection{Two coherent tuples}

All three endpoints use four named coefficient ports,
\((\beta,K_0,M_C,M_F)\), but they use different entire tuples.  The compact
positive tuple is source~44's complement-packed construction
\cite{source44}; the unrestricted-integer tuple is source~34's
mixed-radix, parity-separated construction \cite{source34}.
The binding is summarized below.

\begin{table}[htbp]
\centering
\small
\begin{tabularx}{\textwidth}{@{}lrrX@{}}
\toprule
 & \(P_{13},P_{14}\) & \(Q_{34}\) & Meaning\\
\midrule
Native digit width & 67 & 121 & \(\rho=2^{\text{width}}\)\\
Native cell length & 8,807 & 8,801 & Number of radix-\(\rho\) positions\\
Binary cell width \(d\) & 590,069 & 1,064,921 & \(B=2^d\), \(\beta=B-1\)\\
Permitted native positions & 986 & 989 & Selectors, copies, anchors, dummies\\
Four raw numeral files & 295,388 bytes & 532,479 bytes & Before gzip packaging\\
\bottomrule
\end{tabularx}
\caption{Complete coefficient bindings.  These fixed data are inherited
from the exact recipes and payloads in \cite{source44,source34}.}
\label{tab:coefficients}
\end{table}

In both recipes, if \(D_C,D_R\) denote the native current and shifted
transport coefficients and \(M_{F,\mathrm{native}}\) the native high mask,
the polynomial ports are
\begin{equation}
 \beta=B-1,\qquad K_0=D_C+BD_R,\qquad
 M_C,\qquad M_F=M_{F,\mathrm{native}}+B-1.
 \label{eq:four-port-binding}
\end{equation}
Thus the delivered \(M_F\) is already shifted.  Replacing it by the
unshifted mask, shifting it again, or taking individual ports from the
other tuple changes the polynomial and invalidates the stated binding.
The binding in \path{positive14/ACTIVE_BINDING.json} selects the same
four numeral payloads as \(P_{13}\); the added product register changes
no coefficient.

\subsection{Finite recipes, rather than an unspecified program parameter}

Both recipes start from the same ordered list of \(943\) Wang tiles.
The list is fixed; it is not generated by waiting for a machine to halt.
The compiler sources also relate that literal list to the \(430\)
literal transition actions of the chosen blank machine
\cite{source34,source44,fixed-source}.

For the compact recipe, a tile with ordered labels \((w,n,e,s)\) has
payload
\[
 v_i=w+2^8n+2^{17}e+2^{25}s+127\cdot2^{34},
 \qquad c_i=1+2^{10}(2^{41}-v_i).
\]
Copy and anchor positions supply the \(41\) individual payload bits.
The low ten-bit test forces either no selector or exactly one selector.
Conditional on one selector, a complement sum forces its payload to
equal the copied label and anchor vector.  This is the local clause
mechanism which permits the smaller digit width \cite{source44}.

If \(E\) is that recipe's fixed support and \(c_e\) its specified weights,
the current coefficient and low mask are
\begin{align*}
 D_C={}&\rho^{30}+\rho^{62}+\rho^{1100}+\rho^{2201}
       +\sum_{e\in E}c_e
          \bigl(\rho^{4393-e}+\rho^{5495-e}\bigr),\\
 D_R={}&\rho^{42},\\
 M_C={}&B-1-\sum_{\substack{e\in E\\e\ne1}}\rho^e
                 -2\rho^{1063}.
\end{align*}
Its high mask uses the clause mask \(2^{51}-2\), the specified edge and
anchor tests, and a finite upper padding mask.  The retained source
recipe lists every copy position and every padding exponent.  The
coefficient checkers reconstruct these fixed integers and compare
their raw bytes; the displayed equation does not quantify over
possible choices of \(E\) or of the weights.

The source-34 recipe instead uses parity-separated copy positions
\[
 e_t(j)=o_t+2j,\qquad
 (o_0,o_1,o_2,o_3)=(1020,1000,980,960),
 \qquad 0\le j<10,
\]
and radix-four payload weights \(W_j=2^{10+2j}\).  The four edges are
encoded by their occupied words, and four additional anchor fields
are included.  Its current shifts are \(40,41,1100,2200\), its
transport coefficient is \(D_R=\rho\), and its complete low and high
masks are the ones in the source-34 recipe \cite{source34}.
The different layout supports the ordinary source's integer
normalizations; no assertion that the two layouts are interchangeable
is needed.

\subsection{What the coefficient adapter supplies}

A coefficient recipe alone is not a proof that polynomial roots encode
computations.  The inherited adapters establish the native carry
bounds, the local tile clauses, whole-cell alignment, edge matching,
and the relation between the actual packed rank and the canonical
Pell construction.  Their completeness arguments use finite padding
and a rank controller to arrange the required congruences
\cite{source34,source44,pell-audit}.

For the positive reduction, the canonical compact source in particular
supplies
\[
 q=B^N,\qquad
 J=1+B+\cdots+B^{N-1},\qquad
 N\ \hbox{odd},\qquad R\equiv3\pmod4.
\]
Since \(B=2^{590069}\equiv-1\pmod3\), the odd geometric sum gives
\(J\equiv1\pmod3\).  This exact canonical congruence is what allows the
monic carrier of Lemma~\ref{lem:monic-carrier} to replace the supplied
coordinate \(J\).  The arbitrary-zero soundness proof subsequently
recovers the required integral \(J\); it does not assume the
canonical representation \(q=B^N\) for an arbitrary input.

For the integer construction, the new step is a polynomial bijection
with a fixed source quartic.  All sign and native completeness
obligations of that source are inherited intact; the new proof does
not impose an extra restriction on its zeros.

\section{Signed auxiliary-index recovery}
\label{sec:signed-index}

This section gives the coefficient-independent lemmas used by the
positive construction.  The arguments follow the pinned signed-index
note in package~61 and its explicitly identified Pell sources
\cite{signed-index,pell-audit}.  We include the proofs because the order
of index recovery, rather than just the final norm identities, is
material to the reduction.

\subsection{Pell notation and the elementary dependencies}

For an integer \(A\ge2\), put \(\Delta=A^2-1\) and define
\(\chi_A(n),\psi_A(n)\) by
\[
 \chi_A(n)+\psi_A(n)\sqrt{\Delta}
       =(A+\sqrt{\Delta})^n\qquad(n\ge0).
\]
Equivalently each sequence satisfies
\[
 z(n+2)=2Az(n+1)-z(n),
\]
with \(\chi_A(0)=1,\chi_A(1)=A\) and
\(\psi_A(0)=0,\psi_A(1)=1\).  All these values are integers.  The
addition formulas give
\begin{align*}
 \chi_A(r+s)&=\chi_A(r)\chi_A(s)
                 +\Delta\psi_A(r)\psi_A(s),\\
 \psi_A(r+s)&=\psi_A(r)\chi_A(s)
                 +\chi_A(r)\psi_A(s).
\end{align*}
We use positive integer-coordinate Pell classification and strong
divisibility,
\begin{equation}
 \gcd\bigl(\psi_A(r),\psi_A(s)\bigr)
       =\psi_A\bigl(\gcd(r,s)\bigr).
 \label{eq:strong-divisibility}
\end{equation}
Here \((A,1)\) is the fundamental positive integer-coordinate solution
of \(x^2-\Delta y^2=1\).  The inherited Pell audit collects precisely
these classical ingredients \cite{pell-audit}.

\begin{lemma}[Signed abscissa step-down]
\label{lem:stepdown}
Suppose \(A\ge2\), \(0<r<m\), and \(t\ge0\).  If
\[
 \chi_A(t)\equiv\chi_A(r)\pmod{\chi_A(m)},
\]
then \(t\equiv r\) or \(t\equiv-r\pmod{2m}\).
\end{lemma}

\begin{proof}
Write \(F=\chi_A(m)\).  The addition formulas, together with
\(\chi_A(2m)=2F^2-1\) and
\(\psi_A(2m)=2F\psi_A(m)\), give
\begin{align*}
 \chi_A(j+2m)&\equiv-\chi_A(j)\pmod F,\\
 \chi_A(2m-j)&\equiv-\chi_A(j)\pmod F.
\end{align*}
Reduce \(t\) modulo \(2m\), and reflect an index above \(m\).  Its
abscissa residue has the form
\(\epsilon\chi_A(s)\), where \(0\le s\le m\) and
\(\epsilon\in\{-1,1\}\).
The case \(s=m\) has residue zero, whereas
\(0<\chi_A(r)<F\).

If \(s<m\) and \(\epsilon=-1\), the congruence would force
\(F\mid\chi_A(s)+\chi_A(r)\).  But \(m\ge2\), and the recurrence gives
\[
 0<\chi_A(s)+\chi_A(r)
       \le2\chi_A(m-1)<\chi_A(m)=F.
\]
Thus \(\epsilon=1\).  Equality modulo \(F\) of two positive integers
strictly below \(F\) is equality, and strict growth gives \(s=r\).
The index reductions change \(t\) only by a multiple of \(2m\) and by
sign, proving the assertion.
\end{proof}

\subsection{The odd quotient polynomial}

For \(v\ge0\), define the integer polynomial \(Q_v\) by
\[
 \chi_Z(2v+1)=ZQ_v(Z^2).
\]
The relevant formulas are
\begin{align}
 Q_0(x)&=1,\qquad Q_1(x)=4x-3, \notag\\
 Q_{v+1}(x)&=(4x-2)Q_v(x)-Q_{v-1}(x),\label{eq:Q-recurrence}\\
 Q_v(0)&=(-1)^v(2v+1),\label{eq:Q-zero}\\
 Q_v(1-A^2)&=(-1)^v\psi_A(2v+1).\label{eq:Q-Pell}
\end{align}
The recurrence follows from the Pell addition formulas.  Substitution
of \(x=0\) proves \eqref{eq:Q-zero} by induction.  Substitution of
\(x=1-A^2\), with the odd-index recurrence for \(\psi_A\), proves
\eqref{eq:Q-Pell}; the initial cases are \(1\) and \(1-4A^2\).
These identities are polynomial identities over \(\Z\), so they can
be reduced modulo any positive integer.

\begin{lemma}[Signed auxiliary-index interface]
\label{lem:signed-aux}
Let \(A\ge2\), \(\Delta=A^2-1\), and let
\(p,m,R,c,f,M\) be positive integers satisfying
\begin{gather*}
 c=\psi_A(p),\qquad f=\chi_A(m),\qquad c\mid m,\qquad m>2p,\\
 M\ge2,\qquad c\mid M,\qquad M^2=\Delta(f^2-1),\\
 0<p<c/2,\qquad 0<R<c/2.
\end{gather*}
If \(V,y\in\Z\) satisfy
\begin{equation}
 M^2V^2-(M^2-1)y^2=1,\qquad
 V\equiv-c\pmod f,\qquad V\equiv-R\pmod c,
 \label{eq:aux-interface}
\end{equation}
then \(p=R\).  The signs of \(V\) and \(y\) are unrestricted.
\end{lemma}

\begin{proof}
The norm equation excludes \(V=0\) and \(y=0\).  Apply positive Pell
classification at parameter \(M\) to the positive coordinates
\((M|V|,|y|)\).  There is an index \(\ell\ge1\) with
\[
 M|V|=\chi_M(\ell),\qquad |y|=\psi_M(\ell).
\]
For even indices, the recurrence gives
\(\chi_M(2v)\equiv(-1)^v\pmod M\), which is not zero since \(M\ge2\).
Thus \(\ell=2v+1\) is odd and
\(|V|=Q_v(M^2)\).

As \(M^2\equiv-\Delta=1-A^2\pmod f\), equations
\eqref{eq:Q-Pell} and \eqref{eq:aux-interface} imply
\[
 \psi_A(\ell)\equiv c\quad\hbox{or}\quad-c\pmod f.
\]
Squaring and using
\(\chi_A(2j)=1+2\Delta\psi_A(j)^2\), we obtain
\[
 \chi_A(2\ell)\equiv\chi_A(2p)\pmod{\chi_A(m)}.
\]
Lemma~\ref{lem:stepdown} applies because \(0<2p<m\).
It gives \(2\ell\equiv\pm2p\pmod{2m}\), hence
\(\ell\equiv\pm p\pmod m\), and therefore modulo \(c\).

On the other hand \(c\mid M\), so \(M^2\equiv0\pmod c\).
Equations \eqref{eq:Q-zero} and \eqref{eq:aux-interface} give
\(\ell\equiv\pm R\pmod c\).  Combining the two conclusions yields
\(p\equiv R\) or \(p\equiv-R\pmod c\).
In the first case \(|p-R|<c\), so \(p=R\).
The second case would make \(p+R\) a positive multiple of \(c\),
although \(p+R<c\).  This is impossible.
\end{proof}

The lemma uses only its displayed hypotheses.  It does not require
positivity of \(V\), a previously recovered transport quotient, or
the main-root congruence used later in the positive construction.

\begin{lemma}[Normalized strong norm supplies the divisibility]
\label{lem:normalized-strong}
Suppose \(A\ge2\), \(c=\psi_A(p)>2\), and \(f,I>0\) satisfy
\[
 f^2-(A^2-1)I^2=1,\qquad c^2\mid I.
\]
Then for some \(m\ge1\),
\[
 f=\chi_A(m),\quad I=\psi_A(m),\quad pc\mid m.
\]
In particular \(c\mid m\) and \(m>2p\).  Setting
\(M=(A^2-1)I\) gives \(c\mid M\) and
\(M^2=(A^2-1)(f^2-1)\).
\end{lemma}

\begin{proof}
Pell classification supplies \(m\).
By \eqref{eq:strong-divisibility} and \(c\mid I\), strict growth gives
\(p\mid m\); write \(m=pv\).  With \(D=\chi_A(p)\), composition gives
\[
 I=c\psi_D(v).
\]
Since \(D^2\equiv1\pmod c\), induction in the recurrence yields
\(\psi_D(v)\equiv vD^{v-1}\pmod c\) for \(v\ge1\).
The integer \(D\) is invertible modulo \(c\), and \(c^2\mid I\);
therefore \(c\mid v\).  This proves \(pc\mid m\), and
\(m\ge pc>2p\).  The claims about \(M\) follow directly from the norm.
\end{proof}

This is the normalized argument used for \(P_{13}\), where \(I=ic^2\).
An older lemma for a different ordinary strong norm has an additional
large-\(c\) hypothesis; that hypothesis is not needed or assumed here
\cite{signed-index,pell-audit}.

\subsection{Squared-modulus completeness}

\begin{lemma}[Explicit positive auxiliary reconstruction]
\label{lem:squared-completeness}
Let \(A\ge2\), \(\Delta=A^2-1\), and \(R\ge3\) with
\(R\equiv3\pmod4\).  Suppose
\[
 c=\psi_A(R),\qquad f,I>0,\qquad
 f^2-\Delta I^2=1,\qquad c\mid I.
\]
Put \(M=\Delta I\).  There exist \(T,u>0\) such that, for
\(V=f^2(cT-R)-c\),
\[
 M^2V^2-(M^2-1)(V+u)^2=1.
\]
\end{lemma}

\begin{proof}
Write \(R=2v+1\).  Then \(v\) is odd.  Define
\[
 V_*=\frac{\chi_M(R)}{M}=Q_v(M^2),
 \qquad y_*=\psi_M(R).
\]
These are positive integers, and the Pell identity gives
\[
 M^2V_*^2-(M^2-1)y_*^2=1.
\]
Now \(M^2=\Delta(f^2-1)\equiv1-A^2\pmod{f^2}\).
Since \(v\) is odd, \eqref{eq:Q-Pell} gives
\[
 V_*\equiv-c\pmod{f^2}.
\]
Also \(c\mid M\), so \eqref{eq:Q-zero} gives
\(V_*\equiv-R\pmod c\).
The strong norm gives \(f^2\equiv1\pmod c\), hence
\(\gcd(c,f^2)=1\).  It follows that
\[
 T=\frac{V_*+c+Rf^2}{cf^2}
\]
is a positive integer: its numerator is divisible both by \(c\) and
by \(f^2\).  Its definition gives \(f^2(cT-R)-c=V_*\).

As \(R\ge3\) and \(M\ge2\),
\(V_*\ge4M^2-3>1\).  Rearranging the norm gives
\[
 y_*^2-V_*^2=\frac{V_*^2-1}{M^2-1}>0.
\]
Thus \(u=y_*-V_*\) is a positive integer and \(V_*+u=y_*\),
which completes the construction.
\end{proof}

The modulus \(f^2\) is the reason that the source quotient may be
replaced by \(fT\) in the new positive equation.  Completeness is proved
by constructing a fresh auxiliary solution with this congruence.
No divisibility of an arbitrary pre-existing auxiliary quotient is
assumed.

\section{Positive source coordinates and a congruence inside a Pell root}
\label{sec:positive-coordinate-carrier}

The new equation has thirteen strictly positive unknowns, exact total degree
$28$, and $98$ live binary arithmetic operations.  The degree-$34$,
$93$-operation equation of package~61 is the comparison point.  The change
is an existential coordinate transformation, not a same-coordinate identity.
The inherited native compiler and Pell-index theorems remain explicitly
identified dependencies.

The central choice is to supply the positive gap $C$ instead of the native
index $J$.  Set $q=F+\alpha+C$; the source bounds now hold as soon as the
fixed congruence on $q$ is established.  That congruence is placed inside
the normalized strong Pell root.  A monic square equation recovers the
integrality of that root, and therefore the congruence.  This permits the
transport quotient's sign to be recovered at the end, removing the
predecessor's degree-two positivity multiplier.

\subsection{Literal integer polynomial}

The positive variables are
\[
 F,\alpha,C,f_0,i,T,X,Y,r,\eta,\zeta,u,\gamma.
\]
Use the whole compact source-44 tuple
$(\beta,K_0,M_C,M_F)$, with $\beta=2^{590069}-1$ and the already shifted
high mask $M_F$.  Put
\[
 B=\beta+1,\qquad b=3\beta,\qquad L=\beta b^2.
\]
These are fixed numerals.  The following names are polynomial abbreviations,
not additional witnesses:
\begin{align*}
 q&=F+\alpha+C,& k&=\eta+\zeta,&c&=kY+\eta,\\
 E&=XY,&\ell&=EkY,&a&=E+Y,\\
 H&=4a+3,&\Delta&=a^2+H=(a+2)^2-1,&I&=ic^2,\\
 N_f&=bf_0+q-B,&F_2&=N_f^2,&M_0&=q^3(q-1),\\
 S_\beta&=4\beta qY,&M_\beta&=qS_\beta.&&
\end{align*}
The remaining registers are
\begin{align*}
 G&=(q-1)(q-F)+\alpha,\\
 \widehat R&=\beta G(q^2-1)+(M_C+qM_F)(q-1),\\
 B_t&=(qK_0+X)C-qF,\\
 Z&=M_0(\ell+r)+(q-1)Y-q^3B_t,\\
 \widehat D&=M_\beta(X+ac)
       +H\{M_\beta\gamma+S_\beta X-\beta(k-1)+\widehat R\},\\
 \widehat V&=F_2(\beta cT-\widehat R)-Lc,\\
 \widehat Y&=\widehat V+Lu.
\end{align*}
Define four integer residuals
\begin{align}
 A_0&=Z^2-M_0^2\{\ell(\ell+k)+1\},\label{eq:new-first}\\
 A_1&=\widehat D^{\,2}-M_\beta^2(1+\Delta c^2),\label{eq:new-main}\\
 A_a&=b^2\widehat Y^{\,2}
       -\Delta(F_2-b^2)(\widehat Y^{\,2}-\widehat V^{\,2})-b^2L^2,
       \label{eq:new-aux}\\
 A_s&=F_2-b^2(1+\Delta I^2).\label{eq:new-strong}
\end{align}
The equation is
\begin{equation}
 \boxed{\;P_{13}=A_0^2+A_1^2+A_a^2+A_s^2=0.\;}
 \label{eq:new-p13}
\end{equation}
Every coefficient in this definition is an integer.  There are no variable
exponents or unspecified program parameters.

For the proof only, use the rational abbreviations
\begin{align*}
 f&=N_f/b,&J&=(q-1)/\beta,&R&=\widehat R/\beta,\\
 S&=4qY,&M&=qS=4q^2Y,&v&=k-R-1,\\
 D_*&=M(X+ac)+H(M\gamma+SX-v),&&&\\
 V&=f^2(cT-R)-c,&U&=\Delta(f^2-1).&&
\end{align*}
Then $\widehat D=\beta D_*$, $\widehat V=LV$, and the residuals have the
exact interpretations
\begin{align*}
 A_s&=b^2(f^2-\Delta I^2-1),\\
 A_1&=\beta^2\{D_*^2-M^2(1+\Delta c^2)\},\\
 A_a&=b^2L^2\{(V+u)^2-Uu(2V+u)-1\}.
\end{align*}
These identities clear fixed denominators; they do not introduce divisions
into the delivered polynomial.

\subsection{Recovery of the native congruence}

Suppose \eqref{eq:new-p13} holds.  All four residuals vanish separately.
As $q\ge3$ and $f_0\ge1$,
\[
 N_f\ge2\beta+2>0.
\]
Equation \eqref{eq:new-strong} makes $(N_f/b)^2=1+\Delta I^2$ an integer.
A rational square root of an integer is integral.  Hence $f=N_f/b$ is a
positive integer and
\[
 q-\beta-1=3\beta j,\qquad J=3j+1
\]
for an integer $j$.  The inequality $q>0$ excludes $j<0$.  Thus $J\ge1$,
$q\ge B\ge16$, and, since $\beta\equiv1\pmod3$,
$q\equiv2\pmod3$.  Also $R$ is now integral with its exact native formula.

All of $F,\alpha,C$ are positive and $q=F+\alpha+C$.  The authentic mask
bounds give
\begin{equation}
 3q<R<q^4-2.\label{eq:new-rbounds}
\end{equation}
For example, $G=q\alpha+(q-1)C$ lies between $q$ and $q^2-q$, while
$0<(M_C+qM_F)J<(q-1)(1+2q)$.  These estimates use no transport quotient.

\subsection{The first carrier and normalized rank recovery}

Equation \eqref{eq:new-first} makes $\tau=Z/M_0$ an integer square root of
$\ell(\ell+k)+1$, possibly with negative sign.  Reducing $M_0\mid Z$
modulo $q^3$ and $q-1$ gives
\[
 q^3\mid Y,\qquad q-1\mid B_t.
\]
Write $s=Y/q^3>0$ and $t=B_t/(q-1)\in\mathbb Z$.  Thus
\[
 \tau=\ell+r+s-t,\qquad \tau^2-\ell(\ell+k)=1.
\]
Neither $t$ nor $\tau$ is assumed positive at this stage.

Similarly, \eqref{eq:new-main} gives an integer
$g=D_*/M$ with $g^2-\Delta c^2=1$.  It is positive: $v<k$, $c>kY$,
$H<5a$, and $MY>5$ imply
\[
 D_*>Mac-Hk>ak(MY-5)>0.
\]
Put $A=a+2$ and $P_0=2XY^2+1$.  Pell classification supplies $n,p\ge1$:
\[
 |\tau|=\chi_{P_0}(n),\quad k=2\psi_{P_0}(n),\qquad
 g=\chi_A(p),\quad c=\psi_A(p).
\]
The usual recurrence gives $P_0>A$; since $c>k$, monotonicity implies
$n<p$.  If $n=1$, then $k=2$, $\eta=\zeta=1$, and $c=2Y+1$, whereas
$p\ge2$ forces $c\ge2A\ge4Y+4$.  Hence $n\ge2$, and
\[
 c>kY\ge4P_0Y>8XY^3\ge8q^9>2R,\qquad c>2p.
\]

The \emph{normalized} strong equation is essential here.  It gives
$f=\chi_A(m)$ and $I=\psi_A(m)$.  Since $I=ic^2$, strong divisibility
gives $p\mid m$.  Write $m=pv_0$.  Composition and reduction modulo $c$
give
\[
 I/c=\psi_g(v_0)\equiv v_0g^{v_0-1}\pmod c.
\]
Because $g^2\equiv1\pmod c$, it follows that $c\mid v_0$, hence
$c\mid m$ and $m\ge pc>2p$.  This argument does not require the
large-$c$ hypothesis of the separate ordinary-strong rank lemma.

Set $M_a=\Delta I$.  Then $M_a^2=U$, $c\mid M_a$, and
\eqref{eq:new-aux} becomes
\[
 M_a^2V^2-(M_a^2-1)(V+u)^2=1.
\]
Both displayed Pell coordinates are nonzero.  The definition of $V$ and
$f^2\equiv1\pmod c$ give
\[
 V\equiv-c\pmod f,\qquad V\equiv-R\pmod c.
\]
All hypotheses of the inherited signed auxiliary-index theorem are now
established, including $0<p,R<c/2$ and $c\mid m$.  That theorem yields
$p=R$.  No deleted quotient or transport sign has been used.

\subsection{The adaptive congruence and transport sign}

The main carrier gives $M\mid H(SX-v)$.  Since $q\mid Y$ and
$q\equiv2\pmod3$, we have $\gcd(q,H)=1$.  Moreover, $H$ is odd and
$\gcd(Y,H)\mid3$, so
\[
 \mu=\frac{S}{\gcd(S,H)}\ge\frac{4qY}{3}>q^4>R+1.
\]
Reduction modulo $S$ gives $\mu\mid v$.  Both $S$ and $M$ divide
$4XY^2$ already, because $q^3\mid Y$; no assumption $q\mid X$ is needed.
The first Pell recurrence gives $k\equiv2n\pmod{4XY^2}$.  Thus
$\mu\mid2n-R-1$.  As $n<p=R$ and $|2n-R-1|<R<\mu$,
\[
 2n=R+1.
\]
Therefore $v=k-2n$ is a positive multiple of $4XY^2$, in particular of
$M$.  The main congruence now gives $q\mid HX$, hence $q\mid X$.
The integers
\[
 w=X/q>0,\qquad h=v/E>0,\qquad
 \gamma_{\rm old}=\gamma+w-v/M
\]
restore the source scale and index quotient.  The last one is positive:
$n\ge2$ implies $c>X$, while $g>(a+1)c$, so
$H\gamma_{\rm old}=g-ac-X>c-X>0$.

After $X=qw$, $B_t=q((K_0+w)C-F)$.  As $q-1\mid B_t$ and
$\gcd(q,q-1)=1$, the number
\[
 z_{\rm old}=\frac{(K_0+w)C+q-F-1}{q-1}
\]
is integral.  Its numerator is positive, since $K_0,w,C>0$ and
$q-F=\alpha+C\ge2$.  Thus $z_{\rm old}>0$ and the authentic transport
equation is restored.  This argument allows $z_{\rm old}-1=0$ and does
not assume $w>q$ for arbitrary solutions.

Finally $T_{\rm old}=fT$ and $y=|V+u|$ restore the remaining positive
source coordinates.  The inherited native compiler extracts a finite
syntactic ZFC refutation.

\subsection{Constructive completeness}

Conversely, start with canonical positive normalized source witnesses for
a refutation.  Supply their $F,\alpha,C$ and put
\[
 j=(J-1)/3,\qquad f_0=f-j.
\]
The first is integral, and $f_0>0$ because
$f\ge a+2>Y\ge q^3>q>j$.  This gives $N_f=bf$ and $A_s=0$.
The canonical odd-index reconstruction from package~61 supplies fresh
$T,u>0$ with
\[
 V_* =\frac{\chi_{M_a}(R)}{M_a},\quad y_* =\psi_{M_a}(R),\quad
 T=\frac{V_*+c+Rf^2}{cf^2},\quad u=y_*-V_*.
\]
Its integrality uses the established squared-modulus congruences and
$R\equiv3\pmod4$.  It gives $A_a=0$.

Let $t=q(z_{\rm old}-1)\ge0$ and set
\[
 r=\tau-\ell-s+t,\qquad
 \gamma=\gamma_{\rm old}+v/M-w.
\]
Both are positive.  For the first, $n\ge2$ gives
$k\ge8XY^2+4>3s$, and
\[
 \tau^2-(\ell+k/3)^2
 =k^2(XY^2/3-1/9)+1>0.
\]
Hence $\tau-\ell>k/3>s$.  For the second,
$v$ is a positive multiple of $4XY^2$, so
$v/M\ge XY/q^2=wY/q\ge wq^2>w$.
Substitution gives $Z=M_0\tau$ and $D_*=Mg$, establishing $A_0=A_1=0$.
This proves the reverse existence implication.  Composition with the
native syntactic calibration and negation gives the claimed equivalence
with $\operatorname{Con}(\mathrm{ZFC})$, by ordinary mathematics formalizable
in ZFC.

\subsection{Exact degree and finite verification}

The cancellations
\begin{align*}
 Z&=M_0\ell+Z_t,\\
 A_0&=2M_0\ell Z_t+Z_t^2-M_0^2(\ell k+1),\\
 (X+ac+H\gamma)^2-(1+\Delta c^2)
 &=2ac(X+H\gamma)+(X+H\gamma)^2-Hc^2-1,\\
 \widehat Y^{\,2}-\widehat V^{\,2}&=Lu(2\widehat V+Lu)
\end{align*}
give residual degrees $(13,13,13,14)$.  The degree-$14$ homogeneous part
of $A_s$ is
\[
 -9\beta^2i^2(\eta+\zeta)^4X^2Y^6.
\]
Only $A_s^2$ reaches degree $28$, so the leading form of $P_{13}$ is
\[
 \boxed{81\beta^4i^4(\eta+\zeta)^8X^4Y^{12}.}
\]
The coefficient of $i^4\eta^8X^4Y^{12}$ is $81\beta^4>0$ for the actual
fixed instance.  The exact degree is therefore unaffected by specialization
of the other three native ports.

The independent standard-library checker expands the four literal
residuals into respectively $1921$, $5565$, $26877$, and $81$ symbolic
monomials, compares separately specified formulas, checks their exact
witness degrees, and verifies the final sum of squares over independent
residual indeterminates.  It also checks all wires and liveness.  The
counts are $25$ additions, $16$ subtractions, and $57$ multiplications:
$98$ operations, or $91$ for the four-residual system.  The normal and
optimized Python runs both pass.  These finite checks certify the emitted
algebra and counts; they neither decide consistency nor replace the
unbounded mathematical argument and its named inherited dependencies.


\section{Fourteen positive unknowns in degree twenty-four}
\label{sec:positive14-degree24}

An elementary register lift of $P_{13}$ gives a useful third endpoint:
fourteen strictly positive unknowns, exact degree $24$, and $101$ live
binary arithmetic operations.  Retain all thirteen variables from
Section~\ref{sec:positive-coordinate-carrier} and introduce a positive
variable $E$.  In that section's formulas, use this supplied $E$ wherever
the abbreviation $E=XY$ occurred.  Thus
\[
 a=E+Y,\qquad \ell=E(\eta+\zeta)Y,
 \qquad \Delta=(E+Y)^2+4(E+Y)+3.
\]
Let $\widetilde A_0,\widetilde A_1,\widetilde A_a,\widetilde A_s$ be the
resulting four residuals, and set
\begin{equation}
 \boxed{\;P_{14}=\widetilde A_0^2+\widetilde A_1^2+
   \widetilde A_a^2+\widetilde A_s^2+(E-XY)^2=0.\;}
 \label{eq:positive14-degree24}
\end{equation}
The coefficient tuple is exactly the whole compact tuple already bound to
$P_{13}$.

\paragraph{Exact equivalence.}
Every positive zero $x$ of $P_{13}$ extends to the positive zero $(x,XY)$
of $P_{14}$.  Conversely, a positive zero of $P_{14}$ satisfies $E=XY$,
since every squared residual vanishes separately.  Substitution then
recovers $P_{13}(x)=0$.  The extension is unique, so these maps give a
bijection of positive zero sets.  This elementary argument is formalizable
in ZFC and adds no mathematical source assumption to the preceding
consistency reduction.

\paragraph{Degree and circuit size.}
With $E$ independent, the five residual degrees are respectively
\[
 12,\quad11,\quad12,\quad12,\quad2.
\]
Only $\widetilde A_s$ contains $i$, and its highest homogeneous part is
\[
 -9\beta^2 i^2(\eta+\zeta)^4Y^4(E+Y)^2.
\]
Consequently the coefficient of $i^4\eta^8E^4Y^8$ in $P_{14}$ is
$81\beta^4\ne0$.  The exact total degree is therefore $24$.
The former multiplication $XY$ remains live in the new defining residual.
One subtraction, one square, and one addition extend the $98$-operation
circuit to $101$ operations: $26$ additions, $17$ subtractions, and $58$
multiplications.  The delivered JSON circuit, deterministic generator,
exact symbolic checker, and independent audit make this tradeoff
reproducible.  The register lift is elementary; the substantive new
arithmetic mechanism remains the thirteen-variable construction.


\section{A sextic in thirty-four unrestricted integers}
\label{sec:sextic34}

The second construction improves the inspected degree-six integer equation
from $35$ unknowns and $20$ residuals to $34$ unknowns and $19$ residuals.
Its starting point is the $41$-integer quartic of package~64 at the pinned
repository revision.  The new step is an explicit bijection of its complete
integer zero set, so no new Pell or positivity hypothesis is required.
The inherited quartic's compiler and refutation calibration remain the
semantic dependencies.  No global minimality or new kernel certificate is
claimed.

\subsection{Full polynomial specification}

Use the whole source-34 coefficient tuple $(\beta,K_0,M_C,M_F)$, with
$\beta=2^{1064921}-1$ and the already shifted high mask $M_F$.
This is a different, coherent tuple from the compact tuple used in the
thirteen-positive construction; individual ports are not interchanged.
The $25$ arithmetic unknowns are
\[
 J,F,\alpha,z,f,h,r,v,w,\tau,\eta,\zeta,y,\gamma,o,j,
 E,c,a,D,L,F_2,g,I,u.
\]
There are also three triples of unknowns
$(a_1,a_2,a_3)$, $(b_1,b_2,b_3)$, and $(e_1,e_2,e_3)$.
All $34$ unknowns range over \emph{unrestricted integers}.
Define the abbreviation
\[
 \Pi(x,y,t)=1+x^2+y^2+t(t+1)
\]
and the polynomial registers
\begin{align*}
 q&=\beta J+1,& k&=\eta+\zeta,& C&=q-F-\alpha,& H&=4a+3,\\
 R_L&=\bigl((q-1)(q-F)+\alpha\bigr)L+(M_C+qM_F)J.
\end{align*}
In particular $R_L$ uses the supplied coordinate $L$, not a substituted
degree-two expression for it.  This choice is material to the degree bound.

The first seven residuals are
\begin{align*}
 Q_1&=L-q^2+1,\\
 Q_2&=2u^2-4wq,\\
 Q_3&=a-E-(2v+1)q(L+1),\\
 Q_4&=E(2u^2+1)-2u^2a,\\
 Q_5&=c-\eta-(a-E)k,\\
 Q_6&=\tau^2+(2\tau-k)E(c-\eta)-1,\\
 Q_7&=D-2u^2-\gamma H.
\end{align*}
The next six are
\begin{align*}
 Q_8&=D^2+2acD-Hc^2-1,\\
 Q_9&=k-1-hE-R_L,\\
 Q_{10}&=(K_0+4w)C-F-z(q-1),\\
 Q_{11}&=F_2-f^2,\\
 Q_{12}&=fo-c-jc+k-1-hE,\\
 Q_{13}&=g+F_2o^2-2foc+c^2-y^2.
\end{align*}
Finally put
\begin{align*}
 Q_{14}&=I-(2r+1)c^2,\\
 Q_{15}&=I^2-(a^2+4a+3)(F_2-1),\\
 Q_{16}&=y^2-I^2g-1,\\
 Q_{17}&=F\alpha-\Pi(a_1,a_2,a_3),\\
 Q_{18}&=FC-\Pi(b_1,b_2,b_3),\\
 Q_{19}&=\eta\zeta-\Pi(e_1,e_2,e_3).
\end{align*}
The single equation is
\begin{equation}
 \boxed{\quad Q_{34}=\sum_{t=1}^{19}Q_t^2=0.\quad}
 \label{eq:sextic34}
\end{equation}
The JSON packet calls $v,w$ respectively \texttt{v\_scale},
\texttt{w\_scale}; its three guard triples are the corresponding
\texttt{FA\_pos}, \texttt{FC\_pos}, and \texttt{eta\_zeta\_pos}
coordinates.  The displayed $\tau,D$ are \emph{new} root coordinates,
as the restoration formulas below specify.

\subsection{Reversibility of the reduction}

Set $B_0=E(c-\eta)$ for this proof.  Given a zero of
\eqref{eq:sextic34}, restore the seven deleted source registers by
\begin{align*}
 P&=k+B_0,& X&=2u^2,& N&=q(L+1),\\
 C_2&=c^2,&\Delta&=a^2+4a+3,&V&=fo-c,& B_{\rm aux}&=I^2,
\end{align*}
and restore the source roots by
\[
 \tau_{\rm old}=\tau+B_0,\qquad D_{\rm old}=D+ac.
\]
The source's scale coordinate is $w$ here, so its arithmetic occurrence
of the old scale is $4w$, exactly as in package~64's quartic.
Every restoration is an integer polynomial; no divisibility condition has
been discarded.

The seven defining source rows vanish identically.  The two root shears
give the exact identities
\begin{align*}
 (\tau+B_0)^2-(k+B_0)B_0-1
   &=\tau^2+(2\tau-k)B_0-1,\\
 (D+ac)^2-(a^2+4a+3)c^2-1
   &=D^2+2acD-(4a+3)c^2-1.
\end{align*}
Thus the apparent degree-four terms cancel before the sum of squares is
formed.  The one remaining nontrivial row replacement is
\begin{equation}
 Q_{13}-\bigl((fo-c)^2+g-y^2\bigr)
       =o^2(F_2-f^2)=o^2Q_{11}.
 \label{eq:square-remainder}
\end{equation}
Because $Q_{11}$ is retained, the replaced row and its source row vanish
together.  This is an invertible elementary row operation on the system,
not division by a variable or passage to rational witnesses.  Every other
restored source residual is exactly its listed $Q_t$.

Conversely, given any integer zero of the source quartic, forget its seven
defined registers, set
\[
 \tau=\tau_{\rm old}-E(c-\eta),\qquad D=D_{\rm old}-ac,
\]
and keep all other coordinates.  The same identities establish every
listed residual.  These projection and restoration maps are inverse on
the complete integer zero sets.  Therefore the inherited source theorem
gives, by ordinary mathematics formalizable in ZFC,
\[
 \bigl(\exists\mathbf x\in\mathbb Z^{34}:Q_{34}(\mathbf x)=0\bigr)
 \quad\Longleftrightarrow\quad
 \text{there exists a finite syntactic ZFC refutation}.
\]
Negating yields the required unsolvability/consistency equivalence.

\subsection{Why the degree is exactly six}

All $Q_t$ have degree at most three when the four coefficient ports are
fixed.  In particular retaining $F_2=f^2$ allows the square of $fo-c$ to
be represented by $F_2o^2-2foc+c^2$, which is cubic.  Eliminating $F_2$
instead would lose this saving.  The simultaneous shear of the main root
allows $\Delta$ to be eliminated while retaining that square register.
The construction deletes seven coordinates and seven equations from the
$41$-coordinate, $26$-row source, giving the claimed $34$ and $19$.

Only $Q_{16}=y^2-I^2g-1$ contains both $I$ and $g$.  Hence
\[
 [I^4g^2]Q_{34}=1.
\]
This certifies exact total degree six even after authentic specialization
of all four fixed coefficients.  The expansion with coefficient ports
kept symbolic contains $368$ terms.  No optimized arithmetic-gate claim is
made for this endpoint.

\subsection{Finite verification and proof boundary}

The self-contained standard-library programs verify every source
restoration, the correction \eqref{eq:square-remainder}, all $19$
independently transcribed formulas, liveness of the $34$ variables, and the
unit degree certificate.  They also reconstruct the entire four-numeral
source-34 tuple from the literal $943$-tile list and compare its bytes with
the retained payloads; the four raw integers occupy $532{,}479$ bytes.
Normal and optimized Python runs both pass.

These checks establish exact finite algebra and coefficient integrity.
The unbounded compiler and consistency calibration are the named inherited
mathematical results, and no new end-to-end proof-assistant certificate is
claimed.  The improvement is one unrestricted variable at degree six,
rather than a change in the fewest-variable high-degree bounds.


\section{Inherited theorems and the exact formalization boundary}
\label{sec:proof-status}

\subsection{The ordinary mathematical dependency chain}

There are three distinct levels of justification in this package.
First, the new algebraic reductions are stated and proved in this
report.  The sextic reduction is a zero-set bijection over all integers;
the positive reduction additionally uses classical Pell arithmetic and
the canonical source witness promises explicitly listed above.
Second, the arithmetic source interfaces relate their zeros to finite
native tableaux.  Third, the literal machine/tableau calibration relates
the permitted finite computations to finite syntactic refutations in
the specified ZFC calculus
\cite{source34,source44,source64,fixed-source}.

The inherited components used in the consistency corollaries are:
\begin{enumerate}
 \item \textbf{The complete coefficient adapters.}
 The recipes include carry bounds, exact native supports, local clauses,
 alignment, matching conditions, and canonical rank control.  A numerical
 inequality for each mask does not substitute for this whole interface.
 \item \textbf{The Pell/native arithmetic kernel.}
 Soundness decodes any source zero, while completeness constructs
 suitable positive or signed source coordinates from a finite accepted
 tableau.  The audit in \cite{pell-audit} identifies which Pell lemmas
 are used and how normalized and ordinary strong equations differ.
 \item \textbf{The syntactic refutation calibration.}
 The fixed blank computation is calibrated to the actual finite
 proof relation.  This is the semantic link needed for
 \(\ZFCRef_{\ZFC}\), and is preserved by the new coordinate changes.
\end{enumerate}

The pinned native-calibration note also separates two formal milestones
\cite{native-calibration}.  It records the external theorem
\texttt{KEndpoint215.blank\_halts\_iff\_zfc\_refutation}, relating the
literal blank machine to \texttt{BProv ZFCax\_s [] Form.fBot}.
The formal bridge from the particular coherent native coefficient
instance to that blank-halting predicate remains a distinct contract.
The cited native and Pell sources give ordinary mathematical proofs;
this report does not claim to complete that native Lean contract.

This is a composition of ordinary mathematical proofs.  The report does
not reprove every transitive source theorem, and the exact finite
checkers do not discharge those unbounded dependencies.  The source
manifest identifies the retained files and their pins so that this
boundary is inspectable.  There is no requirement to choose a
coefficient by deciding consistency or by running an unbounded search.

\subsection{The exact \texttt{OBJECT\_ZFC} obligation}
\label{sec:object-language}

ProveIt's pinned \texttt{OBJECT\_ZFC.md} makes a precise distinction
between an external Lean theorem and an internal derivation in its
chosen ZFC calculus \cite{object-zfc}.  In that development, the external
theorem \texttt{ExplicitZFC.noSyntaxZero\_iff\_conZFC} concerns the
actual external predicate
\[
 \texttt{BProv ZFCax\_s [] Form.fBot}.
\]
It is not already an object-language derivation of a sentence expressing
the same equivalence.

The file records the still-open exact target, for that development's
fixed nine-variable polynomial:
\begin{quote}
\small
\begin{verbatim}
BProv ZFCax_s []
  (fIff ObjectCodeSeal.consistencySentence
        ObjectEquation.unsolvabilitySentence).
\end{verbatim}
\end{quote}
Both sides have been supplied as actual \texttt{SetTheory.Form}
sentences, with arithmetic interpretation infrastructure.  The missing
endpoint is the displayed derivation.  Our \(P_{13}\), \(P_{14}\), and \(Q_{34}\)
are different supplied polynomials, so a corresponding certified
endpoint for any of them would also have to instantiate its own equation
sentence and prove the required internal interpretation.

The obstacle cannot be removed by checking standard numerals alone.
An internal ZFC model has its own natural numbers and its own proof
codes, including nonstandard ones.  An external theorem quantified over
the ambient natural numbers does not automatically give correctness
for all those internal objects.  The pinned file identifies the needed
route: establish the internal refutation-checking interpretation and
the arithmetic representation in arbitrary models, or produce
object-language derivations directly, then apply the available
completeness machinery \cite{object-zfc}.

Corollary~\ref{cor:zfc} is the ordinary mathematical claim obtained by
the stated source composition.  This archive supplies no new
end-to-end Lean/Rocq certificate and no exported
\texttt{BProv ZFCax\_s []} proof of that claim.  The new algebra would be
a small and explicit part of such a future formalization; the inherited
compiler and internal proof-coding interpretation remain the larger
obligations.

\subsection{What exact finite verification establishes}

For \(P_{13}\), the delivered circuit and independent symbolic
interpretations establish the four exact residual identities, their
degrees, the fixed-slice leading coefficient, all live inputs and
gates, and the final sum of squares.  The four residual degrees are
\[
 (13,13,13,14),
\]
and their symbolic-port monomial counts are
\[
 (1921,5565,26877,81).
\]
The \(98\) operations consist of \(25\) additions, \(16\) subtractions,
and \(57\) multiplications.  The leading form
\[
 81\beta^4 i^4(\eta+\zeta)^8X^4Y^{12}
\]
shows exact degree \(28\) on the authentic fixed slice, not just a
formal degree upper bound before coefficient specialization.

For \(P_{14}\), the five residual degrees are
\[
 (12,11,12,12,2).
\]
The unique extension \(E=XY\) is checked against the thirteen-coordinate
packet.  The coefficient of \(i^4\eta^8E^4Y^8\) is \(81\beta^4>0\),
which proves exact degree \(24\).  Its \(101\) gates comprise \(26\)
additions, \(17\) subtractions, and \(58\) multiplications; the
five-residual system uses \(92\) gates before its finalizer.

For \(Q_{34}\), exact expansion checks all nineteen independently
specified residuals, the seven restored registers, both root changes,
and the retained-row identity.  The coefficient
\([I^4g^2]Q_{34}=1\) certifies exact degree six independently of the
four fixed ports.  Its symbolic-port expansion has \(368\) terms.
Complete tuple reconstruction and raw-byte comparisons establish the
finite coefficient binding.

The programs are ordinary executable algebra checks.  Their successful
runs do not prove that the polynomial has no zero, exhaust its infinite
domain, or settle \(\Con(\ZFC)\).  They verify the finite formulas and
data to which the unbounded mathematical arguments apply.  The
independent review notes are internal audits with their scope stated;
they are not external peer review.

\section{Other architectures considered}
\label{sec:alternatives}

The construction problem does not require a particular Turing machine
or Jones' universal expression.  We considered routes which could change
the representation more fundamentally.  None produced a stronger
completed bound within this work.  The following distinctions explain
why the proved reductions were retained as the deliverable.

\subsection{A direct proof verifier instead of a blank machine}

A finite ZFC proof checker can be treated as a syntactic computation
without first selecting K215.  One could encode formula parsing,
axiom-instance checking, substitution, and inference directly into
local constraints, or first pass through a small rewriting or
combinatory system.  Such a construction would target the desired
finite-refutation predicate from the beginning.

The size of a verifier description, however, does not directly determine
the number or degree of its final Diophantine unknowns.  A complete
improvement would need an explicit bounded local alphabet, a uniform
packing theorem, sound handling of arbitrarily long proofs, and the
actual fixed coefficients.  Replacing a machine constant while keeping
the same arithmetic skeleton can change coefficient height without
changing the witness count or degree at all.  We did not obtain a
smaller fully calibrated polynomial by this route.  It remains useful
for the internal proof-coding task described in
Section~\ref{sec:object-language}.

\subsection{Universal equations and integer-domain bounds}

Universal Diophantine compilation is a classical baseline
\cite{jones1982}.  Recent work by Bayer, David, Hassler, Matiyasevich,
and Schleicher studies explicit universal bounds over the integers
alongside formalization \cite{bayer2026}; Sun's pinned September 2026
preprint proves undecidability with ten unrestricted integer unknowns
\cite{sun2026}.  These are relevant to the high-degree, few-variable
end of the subject.

The endpoints proved here do not use a new Jones or Sun reduction.
Their arithmetic skeletons still inherit the fixed native/refutation
calibration.  An application of a universal theorem alone would not
establish that its fixed consistency instance improves the degree,
coefficient height, or literal circuit cost of the measured endpoints.
The literature references provide context, not a novelty certificate
for the transformations in this report.

\subsection{Square carriers and free tableau padding}

There is a tempting alternative for integer positivity guards.
Every residue \(a\equiv1\pmod8\) has an odd square root modulo \(2^k\)
for \(k\ge3\), and a chosen root gives arbitrarily large squares in
that residue class.  The elementary lift follows because, for odd \(y\),
\[
 (y+2^{j-1})^2-y^2\equiv2^j\pmod{2^{j+1}}
 \qquad(j\ge3).
\]
Thus, if a computation encoding ignored a sufficiently flexible high
prefix, one might try to move a positivity requirement onto a square
without an additional witness triple.

The existing native word is constrained throughout its masked support.
The source gap \(C\) is a bounded complete word, not merely a low
residue with arbitrary unused high bits.  Arbitrary high digits in a
square can violate local masks, copy fields, and transport relations.
The harmless dummy coordinates used for canonical rank adjustment do
not by themselves prove the freedom needed by this square argument.
No positivity guard was removed on that unsupported premise.

\subsection{Controlled packing collisions}

Lemma~\ref{lem:retained-row} suggests a further direction: a packed
coefficient need only recover a sufficient generating family of
residual equations, rather than isolate every residual literally.
If a collision contributes a polynomial multiple of an equation that
is retained elsewhere, it may preserve the zero set.  This is a
concrete algebraic possibility, as \eqref{eq:retained-square-example}
demonstrates in the sextic.

Extending the idea to a few-variable packing requires additional
proofs about coefficient extraction, signs, and carries at the actual
digit base.  Row-equivalence over a polynomial ring alone does not
authorize a collision in an integer digit encoding.  No improved
nine-variable equation is claimed from that prospective extension.
The sextic gives a completed instance of the retained-equation
principle with all inverse maps explicit.

\section{Reproduction and review}
\label{sec:reproduction}

\subsection{Finite checks and document build}

From the unpacked archive root, run
\begin{quote}
\small
\begin{verbatim}
python3 verify_all.py
python3 -O verify_all.py
\end{verbatim}
\end{quote}
The supplied checkers use the Python standard library.  Source identity,
coefficient reconstruction, literal formula interpretation, exact
degree certificates, and zero-set restoration checks are separate
parts of the replay.  The optimized invocation checks that essential
validation is not silently removed with Python assertions.
The consolidated receipts are \path{VERIFICATION.json} and
\path{VERIFICATION_OPTIMIZED.json}; the unchanged inherited files are
identified in \path{SOURCE_MANIFEST.json}.

The mathematical report is rebuilt by
\begin{quote}
\small
\begin{verbatim}
latexmk -pdf -interaction=nonstopmode -halt-on-error report.tex
\end{verbatim}
\end{quote}
A conventional LaTeX installation with the packages listed in the
preamble is sufficient.  Rebuilding the PDF can change its metadata
and delivery hash without changing the mathematical source.  The
source and coefficient manifests remain the identity records for
those components.

\subsection{Suggested reading order}

The central review question is whether each deletion or changed
coordinate preserves the required integer domain in both directions.
For \(P_{13}\), first inspect recovery of the strong-root congruence,
then the source bounds, then the exact indices, and finally the deleted
quotients and their signs.
In particular, the modulus \(4qY/\gcd(4qY,H)\) is available before
\(q\mid X\); normalized strong divisibility is available before
recovering the old main-root quotient.  The constructive converse
then checks every new positive coordinate.

For \(Q_{34}\), inspect the restored source rows and the inverse root
changes.  The row \(F_2-f^2=0\) is intentionally retained, and the
replacement auxiliary row differs by exactly \(o^2(F_2-f^2)\).
The coefficient of \(I^4g^2\) gives a short exact degree certificate.

After these arguments, compare the literal JSON packets with the
independent formulas and read the coefficient bindings and retained
source proofs.  The archive's root README and source manifest give
the file inventory and provenance.  This order keeps the mathematical
equivalence, its finite representation, and the inherited compilation
boundary separately reviewable.

\begin{thebibliography}{99}
\raggedright
\small

\bibitem{jones1982}
James P. Jones.
\emph{Universal Diophantine equation}.
Journal of Symbolic Logic \textbf{47}(3) (1982), 549--571.
\href{https://doi.org/10.2307/2273588}{doi:10.2307/2273588}.

\bibitem{bayer2026}
Jonas Bayer, Marco David, Malte Hassler, Yuri Matiyasevich, and
Dierk Schleicher.
\emph{Diophantine Equations over \(\mathbb Z\): Universal Bounds and
Parallel Formalization}.
\href{https://arxiv.org/abs/2506.20909v3}{arXiv:2506.20909v3},
22 October 2026.  The version is pinned explicitly.

\bibitem{sun2026}
Zhi-Wei Sun.
\emph{Ten unknowns for Hilbert's tenth problem over the integers}.
\href{https://arxiv.org/abs/2609.01594v1}{arXiv:2609.01594v1},
1 September 2026.  This is the specified first version, not an
unversioned reference.

\bibitem{source34}
ProveIt repository, source 34.
\emph{A mixed-radix, parity-separated, unrounded coefficient tuple}.
\formfile{34-integer14/coefficients/COEFFICIENT_PROOF.md}
{Coefficient proof and complete bounded recipe}.
Repository snapshot \snapshot.

\bibitem{source44}
ProveIt repository, source 44.
\emph{Complement-packed native clauses: a smaller complete coefficient
tuple}.
\formfile{44-b4-sparse/native/PROOF.md}
{Native compiler adapter and coefficient recipe}.
Repository snapshot \snapshot.

\bibitem{pell-audit}
ProveIt repository, source 36.
\formfile{36-small-tuple/proofs/INHERITED_PELL_AUDIT.md}
{\emph{Inherited Pell audit}}.
Repository snapshot \snapshot.

\bibitem{signed-index}
ProveIt repository, source 61.
\formfile{61-congruence-carriers/positive/SIGNED_INDEX_LEMMA.md}
{\emph{Signed auxiliary-index recovery and canonical completeness}}.
The note retains its own earlier source pins; the inspected file is
at repository snapshot \snapshot.

\bibitem{source61}
ProveIt repository, source 61.
\formfile{61-congruence-carriers/README.md}
{\emph{Congruence carriers and a divided Pell norm}}.
Principal thirteen-positive degree-34 equation and its literal
93-gate circuit.  Repository snapshot \snapshot.

\bibitem{source62}
ProveIt repository, source 62.
\formfile{62-sign-and-norm/README.md}
{\emph{Sign recovery and centered norms for explicit ZFC consistency
equations}}.
Thirty-five-integer sextic and thirty-integer octic comparison
artifacts.  Repository snapshot \snapshot.

\bibitem{source64}
ProveIt repository, source 64.
\formfile{64-structural-reductions/README.md}
{\emph{Structural reductions}} and
\formfile{64-structural-reductions/integer41/PROOF41.md}
{the forty-one-integer quartic proof}.
Repository snapshot \snapshot.

\bibitem{source65}
ProveIt repository, source 65.
\formfile{65-shared-taps/README.md}
{\emph{Shared taps}, comparison results}.
Fourteen-positive degree-26, 92-operation endpoint.
Repository snapshot \snapshot.

\bibitem{source58}
ProveIt repository, source 58.
\formfile{58-integral-root/README.md}
{\emph{Integral-root construction}, comparison results}.
Fifteen-positive degree-24, 90-operation endpoint.
Repository snapshot \snapshot.

\bibitem{source60}
ProveIt repository, source 60.
\formfile{60-positive-gap/README.md}
{\emph{Positive-gap construction}, transport companion}.
Fifteen-positive degree-24, 89-operation endpoint; its source-58
parent has 90 operations \cite{source58}.
Repository snapshot \snapshot.

\bibitem{fixed-source}
ProveIt repository.
\formfile{31-sun10/sources/FIXED_BLANK_67_PROOF.md}
{\emph{Fixed blank source proof}}.
Native arithmetic and literal blank-machine finite-window semantics,
with the transitive sources identified there.
Repository snapshot \snapshot.

\bibitem{native-calibration}
ProveIt repository.
\repofile{NATIVE_CALIBRATION52.md}
{\emph{Native arithmetic and the remaining calibration contract}}.
Identifies the external K215/refutation theorem and the remaining
formal native-instance bridge.
Repository snapshot \snapshot.

\bibitem{object-zfc}
ProveIt repository.
\repofile{OBJECT_ZFC.md}{\emph{The object-language ZFC obligation}}.
Exact distinction between the external equivalence and the internal
\texttt{BProv} endpoint, including nonstandard-model obligations.
Repository snapshot \snapshot.

\end{thebibliography}

\end{document}
